\section*{第 \ref{ch:b3:particle-physics} 章——粒子物理}

\begin{solution}{exo:b3:particle-physics:1}
(a) uud：$\tfrac23 + \tfrac23 - \tfrac13 = +1$；udd：$\tfrac23 -
\tfrac13 - \tfrac13 = 0$。(b) $\tfrac23 + \tfrac13 = +1$（
反下\omterm{def:b3:particle-physics:zoo}{夸克}携带 $+\tfrac13$）。(c) $\pi^- =
\bar{\text{u}}\text{d}$；$\Delta^{++} = $ uuu（$3\times\tfrac23 =
+2$——该态的统计行为促使人们引入
“色”这一量子数）。(d) 被拉伸的胶子场储存的能量
$\propto$ 长度；远在\omterm{def:b3:particle-physics:zoo}{夸克}自由之前，储存的能量
超过 $2m_{\text{q}}c^2$，弦便断裂并生成新的
\omterm{def:b3:particle-physics:zoo}{夸克}对：拉出的是\omterm{def:b3:particle-physics:zoo}{介子}，绝不会是裸\omterm{def:b3:particle-physics:zoo}{夸克}。
\end{solution}

\begin{solution}{exo:b3:particle-physics:2}
(a) \omterm{def:b3:particle-physics:zoo}{轻子}：e$^-$、$\nu_\mu$、$\mu^-$；\omterm{def:b3:particle-physics:zoo}{重子}：p、n；\omterm{def:b3:particle-physics:zoo}{介子}：
$\pi^+$。(b) 仅\omterm{def:b3:particle-physics:zoo}{强子}：$\pi^+$、p、n。(c) 单独稳定的：
e$^-$、p 以及中微子（自由中子寿命为十五
分钟；$\mu$ 和 $\pi$ 在微秒量级乃至更短时间内衰变）。(d) 其中
电中性的：$\nu_\mu$ 和 n。
\end{solution}

\begin{solution}{exo:b3:particle-physics:3}
(a) 允许——标准的 $\beta^-$ 衰变（自由中子即如此，
$T_{1/2} \approx \qty{10}{min}$）。(b) 能量上禁止：
产物总质量大于自由质子（在原子核内，结合
能可以抵偿，束缚态的 $\beta^+$ 衰变也会发生）。(c)
禁止：\omterm{prop:b3:particle-physics:conservation}{重子数} $1 \to 0$（以及\omterm{prop:b3:particle-physics:conservation}{轻子数}
$0 \to -1$）。(d) 被\emph{家族}\omterm{prop:b3:particle-physics:conservation}{轻子数}禁止：
μ 子性将消失而电子性出现——从未
被观测到，正因如此而被专门搜寻。
\end{solution}

\begin{solution}{exo:b3:particle-physics:4}
(a) $6\times0.938 \approx \qty{5.6}{GeV}$。(b) 动能
$m_{\text{p}}c^2 \approx \qty{0.94}{GeV}$ 每束
（正面对撞，$\sqrt s = 4m_{\text{p}}c^2$）。(c) 约为 $3:1$：
动量守恒使固定靶中的大部分能量
沦为碎片的无用前向运动——而正面对撞束流则把
能量全部用于碰撞。(d) 1955 年的机器设计为恰好越过
\qty{5.6}{GeV} 的阈值：加速器预算以\omterm{def:b3:relativistic-dynamics:four-momentum}{不变质量}来书写，
而反质子按时出现。
\end{solution}

\begin{solution}{exo:b3:particle-physics:5}
(a) $c\tau = \qty{660}{m}$。(b) \omterm{prop:b3:relativistic-kinematics:dilation}{时间延缓}后的射程
达 ${\sim}\qty{15}{km}$ 需要 $\gamma \sim 23$：$E = \gamma
m_\mu c^2 \sim \qty{2.4}{GeV}$——恰好是宇宙线的能量尺度。(c)
在其自身参考系中，μ 子只存活 \qty{2.2}{\micro s}，
而大气层高速掠过，收缩至
${\sim}\qty{650}{m}$：存活结果相同，只是核算方式不同。(d) \omterm{prop:b3:particle-physics:conservation}{重子
数}：没有足够轻的\omterm{def:b3:particle-physics:zoo}{重子}，而 μ 子是
\omterm{def:b3:particle-physics:zoo}{轻子}；其衰变必须改变味，只有弱相互作用能做到，
因此寿命是从容的微秒量级，而非强相互作用的
幺秒量级。
\end{solution}

\begin{solution}{exo:b3:particle-physics:6}
(a) $\text{d} \to \text{u} + \text{e}^- + \bar\nu_{\text{e}}$：
电荷 $-\tfrac13 = \tfrac23 - 1 + 0$；\omterm{prop:b3:particle-physics:conservation}{重子数}
$\tfrac13 = \tfrac13$；电子家族数 $0 = 1 - 1$。(b) 在两体
衰变中，电子能量已被确定；观测到的
\emph{连续谱}意味着：要么能量守恒失效
（玻尔曾认真考虑过这种可能），或者有第三个不可见的粒子分享了
能量——即泡利提出的中微子，二十五年后被探测到。
(c) $\hbar/m_{\text{W}}c = 197/80400 \approx
\qty{2.5e-3}{fm}$——约为质子的千分之一。(d) 在低
能量下，相互作用必须跨越由其重传递粒子
造成的极短距离，因此振幅极小；在
\qty{100}{GeV} 附近，产生 W 玻色子已成为可能，“弱”相互作用显示出
其真正的电弱强度。
\end{solution}

\begin{solution}{exo:b3:particle-physics:7}
(a) 强相互作用守恒奇异数，因此它只能成对产生
$+1$ 与 $-1$ 的奇异数：一个 K \omterm{def:b3:particle-physics:zoo}{介子}伴随一个超子。(b) 比强相互作用慢了十三
个数量级：唯一破坏味的弱相互作用，
签发了这份死亡证明。
(c) $10^{13}$：一秒对三十万年。
(d) 大量产生表明“强”；而不情愿的衰变则表明“不
强”——只有一个新的守恒数，由强相互作用成对产生，
并由弱相互作用破坏，才能调和二者：这一
记账\emph{就是}这一发现。
\end{solution}

\begin{solution}{exo:b3:particle-physics:8}
(a) 在该对的静止参考系中，总动量为零：单个
光子永远不可能动量为零，而两个背对背的光子可以——
守恒律决定了 PET 的几何构型。(b) $E = mc^2 =
0.002\times(3\times10^8)^2 = \qty{1.8e14}{J}$：两克
物质湮灭所释放的能量超过一千克裂变的铀。(c) 取自
回旋加速器制造的 $\beta^+$ \omterm{def:b3:nuclear-physics:nuclide}{同位素}——$^{18}$F，被组装进
葡萄糖之中：人体自身的化学过程把正电子源置于
耗能最旺盛的细胞中。(d) 这种燃料的成本为其产出的一百万倍，
且无天然矿藏：反物质是探测器的校准工具
兼叙事的引擎，而非能源。
\end{solution}

\begin{solution}{exo:b3:particle-physics:9}
(a) 希格斯势是双阱（即“墨西哥帽”）：
对称点 $\phi = 0$ 是峰顶，真空滚落到
边缘——宇宙的基态打破了自身的
对称性，正如磁体选定北极一样。(b) 这一选择
只打破了部分对称性：残留下来的组合是
电磁相互作用，因此其传递粒子——即光子——保持
零质量，而 W 与 Z，作为被破坏方向的传递粒子，
将相应的劲度系数吸收为质量。(c) \omterm{def:b3:phase-transitions:order}{序参量}的径向
振荡——在磁体中，即
$|m|$ 绕其平衡值的涨落，在
$T_{\text{c}}$ 附近。(d) 宇宙最初的 $10^{-11}$ 秒
内（见\cref{ch:b3:astrophysics}）：高于电弱温度时，
所有粒子都无质量，两种力
原为同一种力。
\end{solution}

\begin{solution}{exo:b3:particle-physics:10}
(a) 每个 \qty{26.7}{MeV} $= \qty{4.28e-12}{J}$ 产生两个
中微子，通量 $= 2\times1360/4.28\times10^{-12} \approx
\qty{6.4e14}{m^{-2}.s^{-1}}$。(b) 以拇指指甲估算：
每秒约 $6\times10^{10}$ 个。一生中穿过你的
${\sim}\qty{0.03}{m^2}$ 面积、厚 \qty{0.3}{m} 的身体，历时
${\sim}\qty{2.5e9}{s}$：约 $5\times10^{22}$，乘以
$10^{-22}\,\unit{m^{-1}}\times\qty{0.3}{m}$：量级为一——
一个太阳中微子一生中就会碰到你一次。(c)
缺失的 $\nu_{\text{e}}$ 在途中已振荡为 $\nu_\mu$ 和
$\nu_\tau$；振荡要求质量有差异，因此
中微子具有质量——这是实验上对
无质量中微子标准模型的首次裂痕。(d) 山体将
一切\emph{其他}粒子全部挡下：只有中微子穿过，因此数公里
岩石把最嘈杂的天空变成了最安静的实验室。
\end{solution}

\begin{solution}{exo:b3:particle-physics:11}
(a) $\gamma = 6.8\times10^{12}/9.38\times10^{8} \approx 7250$：
$1 - v/c \approx 1/2\gamma^2 \approx 10^{-8}$——比光速慢约三米
每秒。(b) $\lambda \approx hc/E \approx
\qty{3e-5}{fm}$：约为质子半径的三万分之一，已进入\omterm{def:b3:particle-physics:zoo}{夸克}
领域。(c) $13.6\,\text{TeV}/938\,\text{MeV} \approx
14500$ 个质子质量——若其全部都能转化。(d) 在
固定能量下 $\gamma_{\text{e}}/\gamma_{\text{p}} =
m_{\text{p}}/m_{\text{e}} \approx 1836$，因此电子损失
高出 $1836^4 \sim 10^{13}$ 倍：质子用于蛮力环形加速器，
电子（干净、点状）则用于较低能量下的精密测量。
\end{solution}

\begin{solution}{exo:b3:particle-physics:12}
(a) $3\times10^{20}\times1.6\times10^{-19} \approx \qty{48}{J}$：
相当于一个被猛力发出的网球——却集中于一个质子之中。(b) $\gamma \approx
3\times10^{11}$；银河系的 $10^{5}$ 光年收缩为
${\sim}10$ 光\emph{秒}：它穿越银河系，按其自身时钟
只需几秒。(c) 在质子参考系中，宇宙微波背景的
温和微波蓝移为能量足够高的 $\gamma$ 射线，
足以从质子上击出 π \omterm{def:b3:particle-physics:zoo}{介子}：高于阈值时宇宙本身
即为吸收体，因此这类质子必然来自近处。(d) 每
平方千米每世纪仅一个此类粒子，种类未知，
无人瞄准：LHC 的 $10^{9}$ 次碰撞每秒，
束流与能量均已知，换来的统计与
控制，是任何宇宙馈赠都无法比拟的。
\end{solution}

\begin{solution}{pb:b3:particle-physics:1}
\textbf{1.} 一个初级质子（通常为 \unit{GeV}--\unit{TeV}）
在高空撞击原子核（强相互作用），溅射出 π \omterm{def:b3:particle-physics:zoo}{介子}；
带电 π \omterm{def:b3:particle-physics:zoo}{介子}经弱相互作用衰变为 μ子与中微子；
μ子仅有轻微电离（电磁相互作用），
向地面疾驰。
\textbf{2.} $c\tau_\pi \approx \qty{8}{m}$：即使 $\gamma$ 很大，
也只能使 π \omterm{def:b3:particle-physics:zoo}{介子}飞行数百米，它们还会在空气中
经强相互作用而消失；μ子不受强相互作用，
寿命长 $85\times$ 倍，因而是能够抵达地面的粒子。
\textbf{3.} $400\,\unit{cm^2}\times1\,\text{min}^{-1} = 400$
每分钟——约每秒 7 个。
\textbf{4.} 噪声闪烁和本地放射性会随机触发
某一块探测器；要使\emph{两块}在 \qty{100}{ns} 之内同时触发，
需要单个粒子穿过这一对探测器——以几何结构充当
一种滤波器。
\textbf{5.} $R_{\text{偶然}} = R_1R_2\Delta t = 100\times100
\times10^{-7} = \qty{e-3}{s^{-1}}$——每一刻钟一次假计数，
相对于每分钟一百个 μ子，可忽略。
\textbf{6.} 只有穿过两块探测器的锥形区域内的
μ子（近乎竖直，约 ${\sim}30$--$40^\circ$ 以内）方可被接受，
这对探测器测量的是方向，而非仅仅是计数率——
这正是它成为\emph{望远镜}的原因。
\textbf{7.} 倾斜的 μ子穿过更厚的大气层：
能量损失更多，飞行中的\omterm{prop:b3:relativistic-kinematics:dilation}{固有时}间更长，因而
存活者更少——大致即你倾斜测量所描绘的 $\cos^2\theta$ 规律。
\textbf{8.} $c\tau = \qty{660}{m}$：按经典理论，从
\qty{15}{km} 处产生的 μ子存活的概率为
$\eu^{-15000/660} \approx 10^{-10}$——阁楼每十年
只能计到一个 μ子。
\textbf{9.} $\gamma = 4000/106 \approx 38$。
\textbf{10.} \omterm{prop:b3:relativistic-kinematics:dilation}{时间延缓}后的寿命 $\gamma\tau \approx
\qty{84}{\micro s}$：射程约 ${\sim}\qty{25}{km}$——阁楼
的计数率恰与相对论的预言一致。
\textbf{11.} 在 μ子参考系中，寿命仍为
\qty{2.2}{\micro s}，而大气层被收缩为
$15000/38 \approx \qty{400}{m}$：穿过它仍有富余的时间。
\textbf{12.} $t_{\text{固有}} \approx 15000/(38\times3\times
10^8) = \qty{1.3}{\micro s}$：存活率 $\eu^{-1.3/2.2} \approx
0.55$——约一半抵达，你的计数器就是它们的普查。
\textbf{13.} 经典的山地与海平面 μ子对比实验
（罗西–霍尔，以及用底片记录的弗里施–史密斯实验）：经典
衰变几乎预言零计数；你所测得的每分钟数百
个计数正是相对论的体现，已在屏蔽层之上得到验证。
\textbf{14.} 第二代带电\omterm{def:b3:particle-physics:zoo}{轻子}：
电荷 $-e$，自旋 $\tfrac12$，质量 \qty{106}{MeV}——即
电子的重孪生兄弟。
\textbf{15.} 电荷：$-1 = -1 + 0 + 0$；电子族：
$0 = 1 + (-1) + 0$；μ子族：$1 = 0 + 0 + 1$。所有
账目都平衡。
\textbf{16.} 三个粒子以任意比例分享能量预算，
因此电子能谱是连续的；它在两个中微子
一同与电子反冲时截止，约为
$m_\mu c^2/2 \approx \qty{53}{MeV}$。
\textbf{17.} 这将破坏\omterm{def:b3:particle-physics:zoo}{轻子}族数——那是标准模型在
无人知其缘由的情况下得以平衡的那本账。寻到它
将是中微子振荡在带电部门的近亲：必然是新物理，
因此实验界竭力追寻它，精度达到
$10^{13}$ 分之一。
\textbf{18.} μ子进入探测器产生一次闪烁；数微秒后，
其衰变电子在同一块探测器中产生第二次闪烁。延迟直方图
是一个指数分布，其斜率\emph{恰为}
$\tau = \qty{2.2}{\micro s}$：在桌面上测得的寿命。
\textbf{19.} 玻尔公式都随质量标度：μ子氢原子
轨道缩小 $207$ 倍，至约 ${\sim}\qty{250}{fm}$——
深入电子云之内，以至于 μ子能感受到
质子的大小：μ子原子是质子的尺子。
\textbf{20.} $10^{15}/10^{6} = \qty{e9}{eV}$：每个粒子约
一 GeV——你阁楼里的 μ子只是普通簇射成员。
\textbf{21.} 相对论性 μ子在水中比光跑得更快，发出
切伦科夫蓝光射入水箱的光电倍增管。一个
$10^{20}\,\unit{eV}$ 簇射的足迹跨越许多
平方千米，因此稀疏的水箱对它取样，犹如
水桶接雨滴。
\textbf{22.} 粒子在介质中沉积能量，介质将其转换为
光，光电倍增管将光转换为
电子，电子学电路计数：从阁楼到对撞机的每一台量能器，
都在重复这同一句话。
\textbf{23.} 弯曲径迹的半径给出动量，其弯曲方向给出
电荷的符号；光子、中子和中微子——
不带电——不留径迹，只在
守恒账目中只留下它们的亏欠。
\textbf{24.} 自然界加速单个质子所达的能量，比
LHC 高一千万倍——令人谦卑——但每平方千米
每世纪仅一个：我们的机器以峰值能量换取
每秒 $10^{9}$ 次可控碰撞。
\textbf{25.} 诞生于高空十千米处的质子；
因\omterm{prop:b3:relativistic-kinematics:dilation}{时间延缓}得以安然抵达（$\gamma \approx 38$，约半数存活）；
由相隔一百纳秒的两次闪烁所确认，并
完全借助一张有十二个座位和四种
力的表格得到解释。
\end{solution}
